\documentclass[11pt]{article}

\usepackage[a4paper,margin=1in]{geometry}
\usepackage{amsmath,amssymb,amsthm,mathtools}
\usepackage{enumitem}
\usepackage{hyperref}
\usepackage{booktabs}
\usepackage{xcolor}
\usepackage{listings}

\lstset{
  basicstyle=\small\ttfamily,
  breaklines=true,
  frame=single,
  backgroundcolor=\color{gray!8}
}

\title{\textbf{Can Tropical Geometry Certify Quantisation?}\\
\large A Research Proposal and Experiment Specification}
\author{}
\date{June 2026}

\newtheorem{theorem}{Theorem}
\newtheorem{proposition}{Proposition}
\newtheorem{definition}{Definition}
\newtheorem{remark}{Remark}

\newcommand{\R}{\mathbb{R}}
\newcommand{\Z}{\mathbb{Z}}
\newcommand{\T}{\mathbb{T}}
\newcommand{\C}{\mathcal{C}}
\newcommand{\B}{\mathcal{B}}

\begin{document}

\maketitle

\begin{abstract}
Post-training quantisation (PTQ) converts a floating-point neural network to an
integer representation for efficient hardware deployment. Existing verification tools
($\alpha$,$\beta$-CROWN, Marabou, Venus) certify pointwise output bounds but do not
check whether the \emph{topological structure} of the decision boundary is preserved.
We propose using the Euler characteristic $\chi$ of the network's tropical hypersurface
as a formal certificate for PTQ. We describe a complete, reproducible experiment
pipeline---using \texttt{canonicalpoly} and SageMath---that tests this certificate
empirically on ReLU networks with INT8 and power-of-two quantisation.
\end{abstract}

\section{Motivation: the gap in PTQ verification}

When a network $f_\theta : \R^d \to \R^k$ is quantised to $f_{\hat\theta}$ with
weights rounded to $b$-bit integers, the standard guarantee is:
\begin{equation}
\sup_{x \in \mathcal{X}} \|f_\theta(x) - f_{\hat\theta}(x)\|_\infty \leq \varepsilon.
\label{eq:output-bound}
\end{equation}
This is a \emph{pointwise} guarantee. It says nothing about whether the decision
regions---the connected components and loops of the classification boundary---are
preserved. A quantised network can satisfy~\eqref{eq:output-bound} while:
\begin{itemize}[leftmargin=*]
\item Merging two decision regions into one (loss of a connected component).
\item Introducing a spurious loop in the boundary (gain of a topological cycle).
\item Splitting a simply-connected region (change in the number of holes).
\end{itemize}
These failures are invisible to output-bound checkers. We propose to detect them
via a topological invariant derived from tropical geometry.

\section{Theoretical foundation}

\subsection{ReLU networks as tropical rational maps}

\begin{theorem}[Zhang, Naitzat, Lim 2018]
Every feedforward ReLU network $f : \R^d \to \R$ is a tropical rational map:
\begin{equation}
f(x) = p(x) \oslash q(x)
:= p(x) - q(x)
\quad \text{in the max-plus semiring,}
\label{eq:tropical-rational}
\end{equation}
where $p, q$ are tropical polynomials (finite max-plus combinations of linear forms).
\end{theorem}

In the max-plus semiring $(\R \cup \{-\infty\}, \oplus, \otimes)$:
\begin{equation}
a \oplus b := \max(a,b), \qquad a \otimes b := a + b.
\end{equation}
A tropical polynomial in $x \in \R^d$ has the form
\begin{equation}
p(x) = \bigoplus_{\alpha \in A} c_\alpha \otimes x^\alpha
= \max_{\alpha \in A} \bigl(c_\alpha + \langle \alpha, x \rangle\bigr),
\end{equation}
where $A \subset \Z^d$ is a finite exponent set and $c_\alpha \in \R$.

\subsection{Decision boundaries as tropical hypersurfaces}

\begin{proposition}[Zhang et al.\ 2018, Proposition 6.1]
Let $f = p \oslash q$ be the tropical rational representation of a ReLU network, and
let $c \in \R$ be the decision threshold. The decision boundary
\begin{equation}
\mathcal{B} := \{x \in \R^d : f(x) = c\}
\end{equation}
satisfies
\begin{equation}
\mathcal{B} \subseteq \mathcal{T}\!\left(c \otimes q \oplus p\right)
= \bigl\{x \in \R^d :
\max\{p(x),\, q(x) + c\} \text{ is attained at least twice}\bigr\}.
\label{eq:tropical-hypersurface}
\end{equation}
\end{proposition}

The right-hand side of~\eqref{eq:tropical-hypersurface} is a \emph{tropical
hypersurface}: the locus where the maximum in a tropical polynomial is attained
simultaneously by two or more monomials. For $d = 2$ it is a piecewise-linear graph.

\subsection{Linear regions as vertices of the canonical polyhedral complex}

\begin{definition}[Canonical polyhedral complex]
For a ReLU network $f$ with $n$ neurons, the \emph{canonical polyhedral complex}
$\C(f)$ is the cell decomposition of $\R^d$ induced by the $n$ hyperplanes
$\{x : W^{(\ell)} x + b^{(\ell)} = 0\}$ arising from the ReLU boundaries at each layer.
Each cell of $\C(f)$ is a convex polytope on which $f$ is affine-linear.
\end{definition}

The vertices of $\C(f)$ are the points where the maximum number of ReLU boundaries
intersect. Masden (2022) showed that the face poset of $\C(f)$ is completely
determined by the \emph{sign sequences} $\sigma(x) \in \{-1, 0, 1\}^n$ of the neurons
at each vertex $x$.

\subsection{The certificate: Euler characteristic of the decision boundary}

For the 1D decision boundary graph $\mathcal{B}$ in $\R^2$, the Betti numbers are:
\begin{equation}
\beta_0(\mathcal{B}) = \text{number of connected components},\quad
\beta_1(\mathcal{B}) = \text{number of independent cycles},
\end{equation}
and the Euler characteristic is
\begin{equation}
\chi(\mathcal{B}) = \beta_0(\mathcal{B}) - \beta_1(\mathcal{B}).
\label{eq:euler}
\end{equation}

\begin{definition}[Tropical PTQ certificate]
Let $f_\theta$ (float32) and $f_{\hat\theta}$ (quantised) be two ReLU networks with
decision boundaries $\mathcal{B}_\theta$ and $\mathcal{B}_{\hat\theta}$. We say the
\emph{tropical PTQ certificate holds} if
\begin{equation}
\chi(\mathcal{B}_\theta) = \chi(\mathcal{B}_{\hat\theta}),
\label{eq:certificate}
\end{equation}
and the \emph{strong certificate holds} if additionally
$\beta_0(\mathcal{B}_\theta) = \beta_0(\mathcal{B}_{\hat\theta})$
and $\beta_1(\mathcal{B}_\theta) = \beta_1(\mathcal{B}_{\hat\theta})$.
\end{definition}

The weak certificate~\eqref{eq:certificate} could hold by cancellation (equal gains
and losses in $\beta_0$ and $\beta_1$) even when the strong certificate fails.
The experiment tests both.

\section{Quantisation schemes}

\subsection{Symmetric INT8}

For a weight matrix $W \in \R^{m \times n}$, symmetric per-tensor INT8 quantisation
rounds each weight to the nearest integer multiple of the scale $s$:
\begin{equation}
s = \frac{\|W\|_\infty}{127}, \qquad
\hat{W} = s \cdot \operatorname{clamp}\!\left(\operatorname{round}\!\left(\frac{W}{s}\right),\, -127,\, 127\right).
\label{eq:int8}
\end{equation}
The rounding error satisfies $\|W - \hat{W}\|_\infty \leq s/2$.

\subsection{Power-of-two (PoT) quantisation}

For ASIC and FPGA deployment, weights are often restricted to powers of two so that
multiplication reduces to a bit-shift:
\begin{equation}
s_{\mathrm{PoT}} = 2^{\lfloor \log_2 \|W\|_\infty \rfloor - (b-1)},
\qquad
\hat{W}_{\mathrm{PoT}} = s_{\mathrm{PoT}} \cdot
\operatorname{clamp}\!\left(\operatorname{round}\!\left(\frac{W}{s_{\mathrm{PoT}}}\right),\,
-(2^{b-1}),\, 2^{b-1}-1\right).
\label{eq:pot}
\end{equation}
PoT introduces larger rounding errors for small weights, so we hypothesise that it
will cause more topological change than INT8.

\section{Experiment pipeline}

The pipeline has six steps. Each step is a standalone script; the full sweep is
orchestrated by a shell script.

\subsection{Step 0 — Environment}

\begin{lstlisting}[language=bash]
conda create -n tropical-ptq python=3.10 -y
conda activate tropical-ptq
conda install -c conda-forge sage=10.3 -y   # ~20 min
pip install torch numpy scipy matplotlib sklearn
git clone https://github.com/mmasden/canonicalpoly
cd canonicalpoly && pip install -e . && cd ..
\end{lstlisting}

Smoke test: construct the simplicial complex $\{\{0,1\},\{1,2\},\{0,2\}\}$
(a triangle) and verify $H_1 \cong \Z$.

\subsection{Step 1 — Train baseline network}

Architecture: $\R^2 \xrightarrow{\mathrm{Linear}_{12}} \mathrm{ReLU}
\xrightarrow{\mathrm{Linear}_{12}} \mathrm{ReLU}
\xrightarrow{\mathrm{Linear}_1} \R$.

Dataset: \texttt{make\_moons}$(n=500, \text{noise}=0.1)$. Trained with Adam for 500
epochs. Required test accuracy $\geq 95\%$.

The Montufar et al.\ (2014) upper bound on the number of linear regions is
\begin{equation}
R(f) \leq \prod_{\ell=1}^{L} \left\lfloor \frac{n_\ell}{n_0} \right\rfloor^{n_0}
\binom{n_0}{0} + \cdots
\end{equation}
For $(n_0, n_1, n_2) = (2, 12, 12)$ this gives approximately $1{,}352$. In practice
we observe 100--400 regions per random seed.

\subsection{Step 2 — Quantise}

Apply~\eqref{eq:int8} and~\eqref{eq:pot} to produce $\hat{W}_{\mathrm{int8}}$ and
$\hat{W}_{\mathrm{PoT}}$. Both yield a standard float32 PyTorch model with
rounded weights---required because \texttt{canonicalpoly} operates on float linear
layers, not PyTorch's packed quantised format.

\subsection{Step 3 — Enumerate the canonical polyhedral complex}

\texttt{canonicalpoly} enumerates all vertices of $\C(f)$ and stores their coordinates
$x^{(v)} \in \R^2$ and sign sequences $\sigma^{(v)} \in \{-1,0,1\}^{24}$.
Two vertices are adjacent in the decision boundary graph iff their sign sequences
differ in exactly one coordinate.

\subsection{Step 4 — Compute Betti numbers}

From the vertex adjacency graph, SageMath computes:
\begin{equation}
H_\bullet(\mathcal{B}) \implies (\beta_0, \beta_1) \implies
\chi = \beta_0 - \beta_1.
\end{equation}

For a typical well-fitted moon network we expect $\beta_0 = 1$ (one connected curve)
and $\beta_1 \in \{1, 2, 3\}$ (a few loops where the boundary wraps around the moons).

\subsection{Step 5 — Certificate comparison}

For each of 10 random seeds, compute:
\begin{equation}
\Delta\chi_{\mathrm{int8}} := \chi(\mathcal{B}_\theta) - \chi(\mathcal{B}_{\hat\theta}^{\mathrm{int8}}),
\qquad
\Delta\chi_{\mathrm{PoT}} := \chi(\mathcal{B}_\theta) - \chi(\mathcal{B}_{\hat\theta}^{\mathrm{PoT}}).
\end{equation}
Certificate holds iff $\Delta\chi = 0$. The experiment reports the empirical rate
and correlates $|\Delta\chi|$ with accuracy drop $\Delta\mathrm{acc}$.

\subsection{Step 6 — Plots}

\begin{description}[leftmargin=1.5em]
\item[Plot A.] Decision boundary before/after PTQ on a $500 \times 500$ grid,
  with enumerated vertices overlaid. Qualitative.

\item[Plot B.] Scatter: $\beta_1(f_\theta)$ vs $\beta_1(f_{\hat\theta})$, one point
  per seed, coloured by $\Delta\mathrm{acc}$. Points on the diagonal = certificate
  holds.

\item[Plot C.] Scatter: $|\Delta\chi|$ vs $\Delta\mathrm{acc}$.
  Pearson $r > 0.5$: certificate is useful.
  Pearson $r < 0.1$: $\chi$ is insufficient; negative result.

\item[Plot D.] Repeat B and C for PoT to test the coarser-rounding hypothesis.
\end{description}

\section{Expected outcomes and their meaning}

\begin{center}
\begin{tabular}{p{0.15\linewidth}p{0.15\linewidth}p{0.55\linewidth}}
\toprule
$\chi$ match & $(\beta_0,\beta_1)$ match & Interpretation \\
\midrule
Yes & Yes & Full topological preservation. Strong certificate holds. \\
Yes & No  & $\Delta\beta_0 = \Delta\beta_1 \neq 0$: cancellation. Weak certificate misleading. \\
No  & ---  & PTQ changed topology. Certificate correctly flags failure. \\
\bottomrule
\end{tabular}
\end{center}

\begin{remark}
Both positive and negative outcomes are publishable. If $\chi$ mismatch correlates
with accuracy drop, the certificate is a viable lightweight alternative to SMT-based
verification for small networks. If it does not, the experiment identifies what a
stronger topological certificate would need to capture (e.g.\ persistent homology
rather than singular homology, or a metric rather than a binary pass/fail).
\end{remark}

\section{Extensions}

\begin{enumerate}[leftmargin=*]
\item \textbf{Higher-dimensional input via random slicing.}
  For $x \in \R^d$ with $d > 2$, sample random 2D affine planes
  $\Pi_1, \ldots, \Pi_K \subset \R^d$, compute $\chi(\mathcal{B} \cap \Pi_i)$ for each,
  and use the average as an approximation to the full invariant. This avoids the NP-hard
  full enumeration.

\item \textbf{Persistent homology certificate.}
  Replace the single-threshold $\chi$ with the persistence diagram
  $\mathrm{dgm}(\mathcal{B})$, which tracks how topological features appear and
  disappear as a filter parameter varies. The bottleneck distance
  $d_B(\mathrm{dgm}(f_\theta), \mathrm{dgm}(f_{\hat\theta}))$ is a continuous certificate.

\item \textbf{Hardware deployment after certificate.}
  After the certificate passes, deploy to an INT8 accelerator (Jetson Nano via TensorRT,
  or a custom FPGA datapath). The certificate then guarantees that the hardware-deployed
  model has the same decision boundary topology as the original floating-point model.

\item \textbf{\texttt{tropical-verify} CLI tool.}
  Package Steps 3--5 as a single command:
  \texttt{tropical-verify --model model.pt --quantized model\_int8.pt}.
  Returns a structured JSON report with $(\beta_0, \beta_1, \chi)$ for both models and
  a pass/fail certificate.
\end{enumerate}

\section{Timeline}

\begin{center}
\begin{tabular}{lll}
\toprule
Day & Task & Acceptance criterion \\
\midrule
1   & Environment: conda + SageMath + canonicalpoly & Smoke test passes \\
2   & Understand canonicalpoly output format & $\chi$ computed for 1 trivial network \\
3   & Train baseline, verify $\geq 95\%$ accuracy & Visual boundary check \\
4--5 & Adapt canonicalpoly to quantised weights & Betti numbers for INT8 model \\
6   & Sweep 10 seeds, collect all Betti numbers & Results table populated \\
7--8 & Debug unexpected results & Root cause identified \\
9--10 & Plots and write-up & Draft complete \\
\bottomrule
\end{tabular}
\end{center}

\begin{thebibliography}{99}

\bibitem{zhang2018}
L.\ Zhang, G.\ Naitzat, L.-H.\ Lim.
\newblock Tropical Geometry of Deep Neural Networks.
\newblock \emph{ICML}, 2018. arXiv:1805.07091.

\bibitem{masden2022}
M.\ Masden.
\newblock Algorithmic Determination of the Combinatorial Structure of the Linear
Regions of ReLU Neural Networks.
\newblock \emph{OpenReview}, 2022. \url{https://github.com/mmasden/canonicalpoly}

\bibitem{montufar2014}
G.\ Montufar, R.\ Pascanu, K.\ Cho, Y.\ Bengio.
\newblock On the Number of Linear Regions of Deep Neural Networks.
\newblock \emph{NeurIPS}, 2014.

\bibitem{tropnnc2025}
K.\ Fotopoulos, P.\ Maragos, P.\ Misiakos.
\newblock TropNNC: Structured Neural Network Compression Using Tropical Geometry.
\newblock \emph{Greeks in AI}, 2025. arXiv:2409.03945.

\bibitem{xu2020}
K.\ Xu et al.
\newblock Automatic Perturbation Analysis for Scalable Certified Robustness and
Beyond ($\alpha$,$\beta$-CROWN).
\newblock \emph{NeurIPS}, 2021.

\bibitem{katz2019}
G.\ Katz et al.
\newblock The Marabou Framework for Verification and Analysis of Deep Neural Networks.
\newblock \emph{CAV}, 2019.

\bibitem{izhakian2010}
Z.\ Izhakian, L.\ Rowen.
\newblock Supertropical Algebra.
\newblock \emph{Advances in Mathematics}, 225(4), 2010.

\end{thebibliography}

\end{document}
